\subsection{求解与结果}
参数 $K$ 表示三级食物链的综合承载条件。固定当前参数组并取代表点后，$K\approx 1.8$ 附近对应规则有界振荡，$K\approx 2.7$ 附近对应明显周期振荡；更高 $K$ 下轨道会呈现更不规则的波动，但严格的混沌识别留到问题四通过最大李雅普诺夫指数完成。后期统计量统一在积分后期时间窗 $t\in[90,180]$ 上提取，包括 $X_1$ 的峰谷差 $\Delta X_1$、平均峰间隔 $\bar T$ 和峰间隔变异系数 $\mathrm{CV}_T$，如表 \ref{tab:q3traj} 所示。

统计量限定在积分后期时间窗内计算，以滤除初值诱发的瞬态影响。由此得到的 $\sigma(X_1)$、$\Delta X_1$ 与 $\mathrm{CV}_T$ 才能用于不同 $K$ 之间的轨道比较。

\begin{table}[htbp]
\centering
\caption{问题三代表轨道的积分后期描述量}
\label{tab:q3traj}
\small
\begin{tabularx}{0.95\textwidth}{>{\raggedright\arraybackslash}p{2.9cm} >{\centering\arraybackslash}p{1.2cm} >{\centering\arraybackslash}p{1.4cm} >{\centering\arraybackslash}p{1.4cm} >{\centering\arraybackslash}p{1.3cm} >{\centering\arraybackslash}p{1.5cm} >{\centering\arraybackslash}p{1.4cm}}
\toprule
类型 & $K$ & $\sigma(X_1)$ & $\Delta X_1$ & $N_{\text{peak}}$ & $\bar T$ & $\mathrm{CV}_T$ \\
\midrule
规则有界振荡 & 1.795 & 0.261 & 0.991 & 4 & 20.720 & 0.004 \\
明显周期振荡 & 2.665 & 0.502 & 2.092 & 3 & 34.178 & 0.014 \\
复杂非线性波动 & 4.800 & 0.885 & 4.093 & 13 & 4.298 & 2.746 \\
\bottomrule
\end{tabularx}
\normalsize
\end{table}

随着 $K$ 增大，$X_1$ 的标准差由 0.261 升至 0.885，峰谷差由 0.991 扩大到 4.093，而峰间隔变异系数由 0.004 升至 2.746，说明系统从规则振荡逐步走向更复杂的长期行为。第三行仅作为复杂非线性波动样本展示，不在问题三中单独给出混沌判断；相关判据仍由问题四的李雅普诺夫扫描提供。图 \ref{fig:q3} 给出三类代表轨道。

\begin{figure}[htbp]
\centering
\includegraphics[width=0.96\textwidth]{fig03_q3_hastings_powell.png}
\caption{问题三代表轨道示意}
\label{fig:q3}
\end{figure}

